\begin{table}[htbp]\centering
\begin{threeparttable}
\caption{Table 1. Local and citywide analyses at a glance}
\begin{tabular}{lcc}
\toprule
 & Local ground-monitor analysis & Citywide satellite analysis \\
\midrule
Treated unit & REMMAQ monitoring stations, with Centro as the main treated station & Quito urban center, UCDB identifier 2544 \\
Comparison pool & Other REMMAQ stations reporting the same pollutant & Latin American urban areas ranked by similarity to Quito \\
Geographic unit & Monitoring station & UCDB urban-center polygon \\
Pollutants & PM2.5, CO, NO2, SO2 & AOD, CO, NO2, SO2 \\
Temporal resolution & Weekly averages of weekday peak-hour readings & Weekly satellite retrievals, aggregated to four-week blocks for estimation \\
Pre-treatment period & 52 weeks before the opening week & 100 weeks before the opening week \\
Post-opening windows & Pre-disruption, donut, and full & Pre-disruption, donut, and full \\
Panel size before estimation & PM2.5: 8 stations x 134 weeks; gases: 7 stations x 134 weeks & Pollutant-specific balanced panels, with 227 weekly observations per city \\
Primary role & Main test of localized air-quality effects near the historic center & Citywide benchmark for aggregate pollution changes in Quito \\
\bottomrule
\end{tabular}
\begin{tablenotes}\footnotesize
\item Local PM2.5 covers all eight monitoring stations. CO, NO2, and SO2 cover seven stations because San Antonio does not report those series. The satellite panels are balanced separately by pollutant.
\end{tablenotes}
\end{threeparttable}
\end{table}
